\documentclass[10pt,a4paper]{amsart}
\usepackage[margin=19mm]{geometry}
\usepackage{amsmath,amssymb,mathtools}
\usepackage[scr=rsfs,cal=euler]{mathalfa}
\usepackage{xfrac}
\usepackage[unicode,hidelinks,bookmarks=false]{hyperref}
\newcommand{\ie}{i.e.\ }
\newcommand{\cf}{cf.\ }
\newcommand{\resp}{respectively\ }
\newcommand{\Subsection}{\subsection}
\providecommand{\nobreakdash}{\nobreak\hbox{-}}
\setlength{\emergencystretch}{2em}
\multlinegap0pt
\usepackage{mathtools}
\usepackage{amssymb}
\usepackage[shortlabels]{enumitem}
\usepackage{xcolor}
\usepackage{tikz}
\newtheorem{theorem}{Theorem}
\newcommand{\sA}{\left.\mathcal{A}_Q\right|_\sigma}
\title{Beyond the Laurent phenomenon}
\author{Andrei Zabolotskii}
\date{Source: arXiv:2603.29073v2 (CC BY 4.0)}
\begin{document}
\maketitle

\noindent\emph{Source and licence.} Beyond the Laurent phenomenon. Version arXiv:2603.29073v2 (CC BY 4.0), distributed by arXiv under the Creative Commons Attribution 4.0 International licence, http://creativecommons.org/licenses/by/4.0/, as recorded in the arXiv licence metadata for this exact version and shown on the abstract page https://arxiv.org/abs/2603.29073v2. Full licence notice: \texttt{licenses/arxiv-2603.29073-CC-BY-4.0.txt}.\par\smallskip

\noindent\textit{Editorial scope.} Counted unit: the article's theorem thm:poly:main (source0.tex lines 148--151). The article's complete original proof of it is delivered with it (source0.tex lines 152--167, one \texttt{proof} environment of 16 lines). No mathematics is written, repaired or re-derived: the statement and the proof are the article's own lines, in the article's own notation, and no retained line is substituted or re-flowed.  Retained uncounted as the source-local context the proof needs: lines 97--100.  Nothing was omitted from any retained span, so the package carries no omission record.  No retained line contains a citation, so the wrapper carries no bibliography and no bibliography entry.  No retained line contains a figure environment, an image-inclusion command or drawing code, so no image is delivered and the package declares no asset dependency.  Editorial formatting only, in addition: an AMS \texttt{amsart} preamble that restates the article's own declarations and macros used by the retained lines, namely the environments \texttt{theorem} on separate counters, together with the standard packages those lines need.  The retained environments print in this document as Theorem 1 in order of appearance, whereas the article prints the same items with the numbering of its own full document.  The complete native source of the article as distributed in the arXiv e-print for this exact version is bundled unchanged as \texttt{sources/197490/source0.tex}.
\begin{equation}
\label{eq:mut}
x_j' = \left.\left( \prod_{i\to j}x_i + \prod_{j\to i} x_i \right)\right/x_j,
\end{equation}

\begin{theorem}
\label{thm:poly:main}
In a partially labelled quiver associated with a specialised cluster algebra $\sA$, \\ a vertex is polyamorous if and only if its neighbourhood is a polycule.
\end{theorem}

\begin{proof}
Suppose $x_j$ is an initial cluster variable in $\mathcal{A}_Q$, $X$ is the collection of all other cluster variables in the initial cluster, $x_j'$ is the new cluster variable arising from the mutation of the initial cluster at $j$, and $y$ is some cluster variable in $\mathcal{A}_Q$. $y$ can be written as a Laurent polynomial $y(x_j,X)$. Suppose $\sigma$ is a labelling that makes the neighbourhood of $j$ a polycule. We need to show that $y|_\sigma$ is a polynomial in $x_j$.

Generally,
\begin{equation}
\label{eq:y}
y = \frac{a_nx_j^n + a_{n-1}x_j^{n-1}+\ldots+a_1x_j+a_0}{bx_j^m},
\end{equation}
where $a_n,\ldots,a_0$ are integer polynomials in $X$, $b$ is a monic monomial in $X$, and $m,n\geqslant0$. By the Laurent phenomenon, $y$ must remain a Laurent polynomial when expressed through the variables of the initial cluster mutated once at $j$, i.e.\ as a function of $x_j'$ and $X$. This implies that the coefficient of each power of $x_j'$ in $y$ expressed that way has to be a Laurent polynomial in $X$.

Per Eq.~\eqref{eq:mut}, $x_j'=P/x_j$ where $P$ is a polynomial in $X$ which is not a monomial (unless $j$ is an isolated vertex, but in that case its neighbourhood cannot be made a polycule) and does not have a nontrivial monomial factor (because $Q$ has no oriented 2-cycles). When $y$ is expressed through $x_j'$ and $X$, each term $a_kx_j^k$ in the numerator of the right-hand side of \eqref{eq:y} becomes $a_k(x_j')^{m-k}P^{k-m}/b$, so when $k<m$, necessarily $a_k = \tilde{a}_kP^{m-k}$ for some polynomial $\tilde{a}_k$ in $X$. However, when $\sigma$ makes the neighbourhood of $j$ a polycule, $P|_\sigma=0$. Therefore each term $a_kx_j^k$ in \eqref{eq:y} with $k<m$ turns into $0$ under specialisation $\sigma$, which makes $y|_\sigma$ a polynomial in $x_j$, as desired.

This establishes polynomiality of the specialisation of all cluster variables, and the polynomiality of all elements of $\sA$ follows immediately.

It remains to show that when $\sigma$ does not make the neighbourhood of $j$ a polycule, not all cluster variables are polynomial in $x_j$. Indeed, in that case $P|_\sigma\ne0$, so we have $x_j'|_\sigma=\left(P|_\sigma\right)/x_j$ non-polynomial in $x_j$.
\end{proof}

\end{document}
